% page 539 of the facsimile (image p-538 of the scan)
% block 185.4 x 246.3 mm ; left margin 16.6 mm
\pgc{-5.080mm}{0.000mm}{
\l{\cel{60}531.}
\l{}
\l{\cel{5}sinkopar. (\sou{netrans}.) I. (patol.) Subisar ceso subita, instantala,}
\l{\cel{8}di la kordio-pulsado, qua eventigas falo pro feblesko maxima. -}
\l{\cel{8}II. (\cel{1}\sou{gram}.) Depreno di un litero od un silabo en la korpo di}
\l{\cel{8}vorto. - III. (\sou{muziko}). Artikular (noto) sur tempo febla, e}
\l{\cel{8}durigar lu til sur tempo forta. - DEFIS.}
\l{}
\l{\cel{5}\sou{sinkretismo}. Probo kunfuzar plura doktrini filozofiala, medicinala.}
\l{\cel{8}DEFIRS}
\l{}
\l{\cel{5}\sou{sinkrona}. I. (\sou{matem.)}\cel{1}Dicesas pri la movi qui eventas en la sama}
\l{\cel{8}tempo-parto. - II. (elektr.) (\sou{pri elektro-mashino}). Di qua la}
\l{\cel{8}angulo-rapideso sempre egalesas la pulso-rapideso di la elektro-}
\l{\cel{8}fluo alternanta quan lu recevas o produktas od esas, di olu,}
\l{\cel{8}multoplo o\cel{2}mutoplo-dividanto. - DEFIS.}
\l{}
\l{\cel{5}\sou{sinodo}. (\sou{religio)}.\cel{2}I. Asemblitaro de membri di la klerikaro. -}
\l{\cel{8}II. Asemblitaro de la kleriki di diocezo,\cel{2}konvokita (advokita)}
\l{\cel{8}da la episkopo. - III. Asemblitaro da la pastori protestantista.}
\l{\cel{8}- DEFIRS.}
\l{}
\l{\cel{5}\sou{sinodika}. (\sou{astron.}) Qua relatas la kunjunto di astri. -\cel{1}\sou{Jiro sino}-}
\l{\cel{8}\sou{dika di Luno}\cel{1}: La tempo-quanto quan lu konsumas por ri-esar en}
\l{\cel{8}la sama situeso relate Suno. -\cel{1}\sou{Yaro sinodika}\cel{1}: Qua ri-esigas}
\l{\cel{8}Tero ye la sama longitudo kam altra planeto. - DEFIRS.}
\l{}
\l{\cel{5}\sou{sinonima}. Dicesas pri vorto quan karakterizas lua analogeso gene-}
\l{\cel{8}rala, pri senco, kun altra vorto, ma diferas de ica per nuanco}
\l{\cel{8}di signifiko. - DEFIRS.}
\l{}
\l{\cel{5}\sou{sinoptika}. I. Qua igas videbla, per un regard-uno, omna parti di}
\l{\cel{8}ensemblo. (Kom ex. La tabelo sinoptika di la\cel{1}\sou{Kurso di la filozo}-}
\l{\cel{8}\sou{fio pozitivista}, da Auguste Comte).- II. (\sou{kristanismo}.)}
\l{\cel{8}La\cel{1}\sou{sinoptiki}\cel{1}:La evangelii da Metteus, da Markus, da Lukas e}
\l{\cel{8}da Johannes, qui esas dispozita same en la ensemblo de la naraco.}
\l{\cel{8}- DEFIRS.}
\l{}
\l{\cel{5}\sou{sinovio}. (\sou{anat.}) Humoro di qua la rolo esas igar plu facila la}
\l{\cel{8}interglito di la artiki moviva. - DEFIRS.}
\l{}
\l{\cel{5}\sou{sinso}. De dextre a sinistre, pasante avan la observanto, od}
\l{\cel{8}inverse. -(Sur la mapi, onu kustumas indikar la sinsi per}
\l{\cel{8}flechi). -}
\l{}
\l{\cel{5}\sou{sintaxo.}\cel{2}(\sou{gram.}) Ensemblo ek la reguli qui regnas la dispozo di}
\l{\cel{8}la vorti en la frazo, e la konstrukto di la propozicioni. -}
\l{\cel{8}\sou{Sintaxo di koordino}:Qua regnas la konstrukto di la propozicioni}
\l{\cel{8}inter-juntita.-\cel{1}\sou{Sintaxo di subordino}\cel{1}: Qua regnas la konstrukto}
\l{\cel{8}di la propozicioni inter-dependanta. - DEFIRS.}
\l{}
\l{\cel{5}\sou{sintezar}. (\sou{trans.)}\cel{2}I. Ri-kompozar la elementi di ensemblo qui}
\l{\cel{8}esas separita; (a) (\sou{kemio}) Rikonstruktar korpo kompozura, di qua}
\l{\cel{8}la elementi esas separita per analizo ; (b) (\sou{kirurgio}) Riunionar}
\l{\cel{8}la parti apartigita di vunduro, la fragmenti di osto ruptita ;}
\l{\cel{8}(c) (\sou{logiko}) Komencar per lo generala, de qua onu obtenas lo}
\l{\cel{8}partikulara per nura dedukto logiko-konforma. - DEFIRS.}
}
